% ============================================================================
%  preamble/base.tex  --  the core setup every deck needs
% ----------------------------------------------------------------------------
%  Load this first. It pulls in the engine/font setup, then the small set of
%  packages that essentially every slide deck uses.
%
%  Optional modules (maths, code, figures, bibliography, Thai) are separate
%  files -- load only the ones a given deck actually needs.
% ============================================================================

% Safety net: main.tex normally sets this (see decks/_template/main.tex).
\providecommand{\deckroot}{../../}

% --- Engine + font setup (must come first) ---------------------------------
% ============================================================================
%  preamble/engine.tex  --  make one source file work on all three engines
% ----------------------------------------------------------------------------
%  LaTeX has three modern compilers and they handle fonts completely
%  differently:
%
%    pdfLaTeX   old, fast, uses pre-built Type1 fonts. Cannot read .ttf/.otf
%               files, so it CANNOT use Sarabun.
%    XeLaTeX    can use any .ttf/.otf font on disk. Needed for Sarabun.
%    LuaLaTeX   same font abilities as XeLaTeX, plus scripting.
%
%  This file detects which one is running and sets up the text encoding
%  accordingly, so the same main.tex compiles under any of them with no edits.
% ============================================================================

\usepackage{iftex}

\ifPDFTeX
  % ------------------------------------------------------------------------
  %  pdfLaTeX branch
  %  T1 fontenc gives proper accented characters and correct hyphenation.
  %
  %  NOTE: \usepackage[utf8]{inputenc} is deliberately NOT loaded.
  %  UTF-8 has been LaTeX's default input encoding since the 2018 release,
  %  so that line is a no-op on any current TeX Live -- and it is an error
  %  under XeLaTeX/LuaLaTeX. Leaving it out keeps this file engine-neutral.
  % ------------------------------------------------------------------------
  \usepackage[T1]{fontenc}
\else
  % ------------------------------------------------------------------------
  %  XeLaTeX / LuaLaTeX branch
  %  fontspec is what lets LaTeX open .ttf/.otf files directly.
  % ------------------------------------------------------------------------
  \usepackage{fontspec}
  \defaultfontfeatures{Ligatures=TeX}
\fi

% ---------------------------------------------------------------------------
%  TYPEFACES
%  Fira Sans ships with TeX Live in BOTH Type1 and OpenType form, and these
%  two packages pick the right one for whichever engine is running. That is
%  why the same two lines work everywhere and do not need to be inside the
%  branch above.
% ---------------------------------------------------------------------------
\usepackage[sfdefault]{FiraSans}   % body + headings
\usepackage{FiraMono}              % code; sets \ttdefault by itself

\renewcommand{\familydefault}{\sfdefault}

% ---------------------------------------------------------------------------
%  MATHS
%  Fira has no matching maths font, so maths stays in Computer Modern. That
%  is deliberate: it is what most maths papers look like and it reads well
%  on a projector.
% ---------------------------------------------------------------------------


% --- Text ------------------------------------------------------------------
\usepackage{verbatim}      % verbatim text, and the handy `comment` environment
\usepackage{array}         % better tabular column types: >{} <{} m{} b{}
\usepackage{booktabs}      % \toprule \midrule \bottomrule -- proper table rules
\usepackage{multirow}
\usepackage{ragged2e}      % \justifying, better \raggedright with hyphenation
\usepackage{microtype}     % subtle spacing improvements; free quality win
\usepackage{algorithm2e}   % algorithm
\usepackage{caption}       % caption

% --- Graphics --------------------------------------------------------------
\usepackage{graphicx}
%  Where \includegraphics looks, in order: this deck's figures/, the deck
%  folder itself, then the shared theme assets (so \includegraphics{logo...}
%  works from any deck without spelling out a path).
\graphicspath{{figures/}{./}{\deckroot theme/assets/}}

% --- Slide utilities -------------------------------------------------------
\usepackage{appendixnumberbeamer}  % stops the appendix inflating the slide count

% ---------------------------------------------------------------------------
%  WARNING SUPPRESSION
%  The `silence` package can mute LaTeX warnings. Used with care it keeps the
%  build log readable; used as a blanket mute it hides real problems.
%
%  Each line below is targeted and explained. Do not add \WarningsOff with no
%  argument -- that turns off ALL warnings, including ones you need to see.
% ---------------------------------------------------------------------------
\usepackage{silence}


% -------
\usepackage{colortbl}

% Beamer warns that \RequirePackage'ing a font package in a theme "may not
% work as expected". It does work here; the warning is noise on every run.
\WarningFilter{beamer}{Unknown document class option}

% Beamer's own "Font shape ... undefined" chatter when a bold small-caps
% variant is requested and substituted. The substitution is correct.
\WarningFilter{latexfont}{Font shape}

% ---------------------------------------------------------------------------
%  Nicer defaults
% ---------------------------------------------------------------------------
\setlength{\parskip}{0.5em}
\setlength{\parindent}{0pt}

% ---------------------------------------------------------------------------
%  LINE SPACING
%  \decklinespread{1.15} from your deck tightens every line on every slide.
%  Left alone it is LaTeX's 1.0 -- except that preamble/lang-thai.tex raises
%  it to 1.25, because Thai stacks vowel and tone marks above the baseline and
%  hangs vowels below it, and Latin leading is not tall enough to clear both.
%
%  Below about 1.15 those marks start touching the descenders of the line
%  above, so lower it only on a deck with no Thai in it.
%
%  \selectfont is what makes the change take effect; \linespread on its own
%  is stored and ignored until the next font selection.
% ---------------------------------------------------------------------------
\newcommand{\decklinespread}[1]{\linespread{#1}\selectfont}
